\section{Global soft reassembly of the actual static mask}
\label{sec:tpc36-global-mask}

We now handle the three factors of
\eqref{eq:tpc36-actual-static-mask}.  Every identity in this section is
pointwise on the finite physical row set.

\subsection{Same-source deletion}

Let \(\cL\) be the finite set of prime sources in the packet.  Inject it
into \(\Z/M_L\Z\).  Orthogonality gives
\begin{equation}
 \one_{\ell=\ell'}
 =\frac1{M_L}\sum_{t\bmod M_L}
 e_{M_L}(t\iota(\ell))e_{M_L}(-t\iota(\ell')).
 \label{eq:tpc36-source-equality}
\end{equation}
Hence
\begin{equation}
 \pinf(\one_{\ell=\ell'})=1,
 \qquad
 \pinf(\one_{\ell\ne\ell'})\le2.
 \label{eq:tpc36-source-cost}
\end{equation}
The equality in the first formula follows also from the entrywise lower
bound.  There is no factor equal to the number of prime sources.

\subsection{Near-row deletion}

Order the \(N\) physical rows so that
\begin{equation}
 m_1<m_2<\cdots<m_N,
 \label{eq:tpc36-row-order}
\end{equation}
which is possible by the inherited row-label injectivity.  Write
\(r(\alpha)\) for the rank of \(m_\alpha\).  For a fixed right row
\(\gamma\), define
\begin{align}
 k_-(\gamma)
 &:=\#\{\alpha:m_\alpha<m_\gamma-\Delta\},
 \label{eq:tpc36-lower-rank}\\
 k_+(\gamma)
 &:=\#\{\alpha:m_\alpha\le m_\gamma+\Delta\}.
 \label{eq:tpc36-upper-rank}
\end{align}
Then, including all endpoint conventions exactly,
\begin{equation}
 \boxed{
 \one_{|m_\alpha-m_\gamma|>\Delta}
 =1+H_N(k_-(\gamma),r(\alpha))
    -H_N(k_+(\gamma),r(\alpha)).}
 \label{eq:tpc36-near-row-prefix}
\end{equation}

\begin{proposition}[Near-row Fourier reassembly]
\label{prop:tpc36-near-row-projective}
For every \(\Delta\ge0\),
\begin{equation}
 \pinf\!\left(\one_{|m_\alpha-m_\gamma|>\Delta}\right)
 \ll1+\log(2N)\ll\log(2Q).
 \label{eq:tpc36-near-row-projective}
\end{equation}
\end{proposition}

\begin{proof}
Insert \eqref{eq:tpc36-prefix-Fourier} into the two prefix kernels in
\eqref{eq:tpc36-near-row-prefix}, then apply the triangle inequality.
The physical row count is \(N\ll QX^{o(1)}\), so
\(\log(2N)\ll\log(2Q)\).
\end{proof}

The proof uses only the order of the row labels.  It does not require the
near band to have constant width or the labels to fill an integer interval.

\subsection{Orthogonal compression of the row-gcd projector}

TPC-33 proves the exact identity
\begin{equation}
 \one_{(d,e)\le G}
 =\sum_{\substack{v\mid d\\v\mid e}}\lambda_G(v),
 \qquad
 \lambda_G(v)=\sum_{\substack{g\mid v\\g\le G}}\mu(v/g).
 \label{eq:tpc36-gcd-projector}
\end{equation}
We upgrade its fixed-column absolute complexity to a global signed
projective representation.

Let \(\cI_G\) be the finite set of \(v\) dividing at least one physical
squarefree row coordinate, and inject \(\cI_G\) into
\(\Z/M_G\Z\).  Put
\begin{align}
 P_t(d)
 &:=\sum_{v\mid d}
 \operatorname{sgn}(\lambda_G(v))
 \sqrt{|\lambda_G(v)|}\,e_{M_G}(t\iota(v)),
 \label{eq:tpc36-gcd-P}\\
 Q_t(e)
 &:=\sum_{v\mid e}
 \sqrt{|\lambda_G(v)|}\,e_{M_G}(-t\iota(v)),
 \label{eq:tpc36-gcd-Q}
\end{align}
with \(\operatorname{sgn}(0)=0\).

\begin{theorem}[Global orthogonal gcd compression]
\label{thm:tpc36-global-gcd}
For physical squarefree \(d,e\),
\begin{equation}
 \boxed{
 \one_{(d,e)\le G}
 =\frac1{M_G}\sum_{t\bmod M_G}P_t(d)Q_t(e).}
 \label{eq:tpc36-global-gcd-factorization}
\end{equation}
If
\(\omega_*=\max\{\omega(d):d\text{ occurs in the packet}\}\), then
\begin{equation}
 \boxed{
 \pinf\!\left(\one_{(d,e)\le G}\right)
 \le(3+2\sqrt2)^{\omega_*}=X^{o(1)}.}
 \label{eq:tpc36-global-gcd-cost}
\end{equation}
\end{theorem}

\begin{proof}
Insert \eqref{eq:tpc36-gcd-P}--\eqref{eq:tpc36-gcd-Q}, average over
\(t\), and apply \eqref{eq:tpc36-finite-orthogonality}.  Only equal
divisors survive, giving \(\sum_{v\mid d,e}\lambda_G(v)\), which is
\eqref{eq:tpc36-gcd-projector}.

For squarefree \(v\),
\begin{equation}
 |\lambda_G(v)|
 \le\sum_{g\mid v}|\mu(v/g)|
 =2^{\omega(v)}.
 \label{eq:tpc36-lambda-bound}
\end{equation}
Therefore, uniformly in \(t\),
\begin{align}
 |P_t(d)|
 &\le\sum_{v\mid d}2^{\omega(v)/2}
 =(1+\sqrt2)^{\omega(d)},
 \label{eq:tpc36-P-bound}\\
 |Q_t(e)|
 &\le(1+\sqrt2)^{\omega(e)}.
 \label{eq:tpc36-Q-bound}
\end{align}
The normalized counting measure in
\eqref{eq:tpc36-global-gcd-factorization} has total mass one.  Multiplying
the two suprema gives
\((1+\sqrt2)^{2\omega_*}=(3+2\sqrt2)^{\omega_*}\).
Finally, \(C^{\omega(n)}\ll_{C,\eps}n^\eps\) for each fixed \(C>1\),
and \(d,e\ll X^{O(1)}\).
\end{proof}

\begin{remark}[Why absolute divisor expansion looked expensive]
The raw set \(\cI_G\) can have polynomial size.  Summing one rank-one
term per \(v\) in absolute value ignores orthogonality between divisor
coordinates.  The average in
\eqref{eq:tpc36-global-gcd-factorization} has the same number of formal
phases but normalized mass one; only the divisor suprema remain.
\end{remark}

\subsection{The complete static-mask theorem}

\begin{theorem}[Soft projective reassembly of the actual mask]
\label{thm:tpc36-global-static-mask}
On every inherited physical packet,
\begin{equation}
 \boxed{
 \pinf(\mathfrak m_{\rm act})
 \ll\log(2Q)(3+2\sqrt2)^{\omega_*}
 =X^{o(1)}.}
 \label{eq:tpc36-global-static-mask-bound}
\end{equation}
The estimate is uniform in the fixed pruning parameter and its threshold
\(G=X^{\kappa_0}\).
\end{theorem}

\begin{proof}
The mask \eqref{eq:tpc36-actual-static-mask} is the Schur product of the
three kernels treated in \eqref{eq:tpc36-source-cost},
\eqref{eq:tpc36-near-row-projective}, and
\eqref{eq:tpc36-global-gcd-cost}.  Apply
\eqref{eq:tpc36-schur-submultiplicativity}.
\end{proof}

\begin{corollary}[The retained-multiplier interface]
\label{cor:tpc36-full-multiplier-reassembly}
After multiplication by the actual fixed-period and smooth factors in
\eqref{eq:tpc36-actual-joint-multiplier}, the joint multiplier has the
residue-cell representation
\begin{equation}
 \mathfrak A_{\alpha,\gamma}^{\rm act}(j)
 =\sum_{r\bmod H_0}\one_{j\equiv r\!\!\!\pmod{H_0}}
 \int_{\Omega_r}p_{\omega,r}(\alpha)q_{\omega,r}(\gamma)
 \phi_{\omega,r}(j/J)\,d\nu_r(\omega)
 \label{eq:tpc36-full-multiplier-reassembly}
\end{equation}
such that, for every fixed \(B\ge0\),
\begin{equation}
 \sum_{r\bmod H_0}\int_{\Omega_r}
 \|p_{\omega,r}\|_\infty\|q_{\omega,r}\|_\infty
 \|\phi_{\omega,r}\|_{C^B}\,d|\nu_r|(\omega)
 \ll_B X^{o(1)}.
 \label{eq:tpc36-orbit-seminorm-cost}
\end{equation}
Thus, after the fixed \(O_h(1)\) residue-cell split, the conditional
projective product representation of TPC-32 is unconditional for the
actual multiplier on this packet.
\end{corollary}

\begin{proof}
TPC-25 supplies the finite residue-cell split and, on each cell, an
absolutely convergent Mellin separation for the smooth row--orbit factors,
with every fixed orbit seminorm integrable \citep{WangTPC25}.  The
periodic \(j\)-factor is constant on a cell and is not differentiated as a
function of \(j/J\).  Take the product with the representation in
\cref{thm:tpc36-global-static-mask}; the static factors do not change the
orbit seminorms.  Summing the fixed number of cells proves the claim.
\end{proof}

\begin{corollary}[A trace-norm consequence]
\label{cor:tpc36-static-trace-norm}
If the physical row set has cardinality \(N\), then
\begin{equation}
 \|\mathfrak m_{\rm act}\|_*
 \le N\pinf(\mathfrak m_{\rm act})
 \ll QX^{o(1)}.
 \label{eq:tpc36-static-trace-bound}
\end{equation}
\end{corollary}

This is an upper bound for the static mask matrix.  It is not an upper
bound of size \(Q^2/J\) for either orbit-Gram off-diagonal matrix, whose
entries contain an additional sum over time and arithmetic coefficients.
